\documentclass[12pt]{article}
\usepackage[utf8]{inputenc}
\usepackage[T1]{fontenc}
\usepackage{hyperref}
\title{Yimchungru Naga Language Resources for K-12 Curriculum Development}
\author{Compiled by Claude Sonnet 5.5 (claude-sonnet-5-5) for Andr\'as Kornai}
\date{}
\begin{document}
\maketitle

\section{Introduction}
Yimchungru Naga (also Yimchungrü, Yimkhiung; ISO 639-3 yim; Glottocode yimc1240) is a Sino-Tibetan language of Nagaland, India. Speaker estimates in the sources conflict: the task description gives about 92,000, Joshua Project about 63,000, and Global Recordings Network 37,000 for the Wai variety, and these figures measure different things. The search found a complete Bible and some audio, but no school textbooks or primers.

\section{Primary Educational Resources}
No primers, textbooks or school programmes were found in the search results.

\section{Linguistic and Reference Materials}
\begin{itemize}
\item B.\,B. Kumar, \emph{Hindi Yimchungr\"u English dictionary} (Kohima: Nagaland Bhasha Parishad, 1973), with 363 records in the STEDT database. \url{https://stedt.berkeley.edu/~stedt-cgi/rootcanal.pl/source/BBK-Yim}
\item Endangered Languages Project, Yimchungru Naga: a 38-word vocabulary and a link to a 1967 classification document with wordlists. \url{https://endangeredlanguages.com/lang/5709/guide/9489}
\end{itemize}

\section{Religious and Community Resources}
\begin{itemize}
\item Bible Society of India and the Yimchungr\"u Baptist Church Association completed a full Bible translation (Yimkhiung Naga Bible, 2007). A Bible app listing appeared in the results; I did not confirm which text it carries. \url{https://www.bible.com/ca/versions/1063}
\item Joshua Project profile of the Yimchungra in India. \url{https://joshuaproject.net/people_groups/18347}
\end{itemize}

\section{Digital and Audio Resources}
\begin{itemize}
\item Global Recordings Network, Yimchungru Naga (Wai variety): audio Bible stories (Words of Life). \url{https://globalrecordings.net/en/language/14622}
\item Global Recordings Network, Yimchungru Naga (Pherrongre variety). \url{https://www.globalrecordings.net/en/language/14620}
\end{itemize}

\section{Curriculum Development Context}
The language is vigorous in its community, but the available published material is limited to a 1973 dictionary, scripture, and mission audio. A school curriculum would need new primers, a standardized orthography decision, and local writers. Dialect varieties (Wai, Pherrongre) should be kept distinct. The Baptist church association is a likely community partner, and the 2007 Bible is the largest text in the language I found.
\end{document}